% !TEX program = xelatex
\documentclass[12pt,a4paper,oneside]{report}

% =========================================================
% 1. CẤU HÌNH HỆ THỐNG & NGÔN NGỮ
% =========================================================
\usepackage{fontspec}
\usepackage[vietnamese]{babel}
\setmainfont{Times New Roman}

\usepackage[utf8]{inputenc}
\usepackage{amsmath, amssymb, amsfonts, amsthm}
\usepackage{geometry}
\geometry{a4paper, margin=2.5cm, headheight=15pt}
\usepackage{setspace}
\onehalfspacing % Giãn dòng 1.5 chuẩn luận văn

% =========================================================
% 2. GÓI BỔ TRỢ ĐỒ HỌA & TRÌNH BÀY
% =========================================================
\usepackage{graphicx}
\usepackage{booktabs}
\usepackage{longtable}
\usepackage{caption}
\usepackage{subcaption}
\usepackage{float}
\usepackage{xcolor}
\usepackage{hyperref}
\usepackage[nameinlink,capitalize]{cleveref} % Cross-referencing thông minh
\usepackage[most]{tcolorbox} % Khung hộp cho Case Study

% Định dạng màu sắc
\definecolor{matchaiblue}{RGB}{0, 51, 102}
\definecolor{casegrey}{RGB}{245, 245, 245}

% Cấu hình Hyperref
\hypersetup{
    colorlinks=true,
    linkcolor=matchaiblue,
    filecolor=magenta,      
    urlcolor=cyan,
    pdftitle={MatchAI Pro Monograph},
}

% Cấu hình TColorBox cho Case Study
\newtcolorbox{casebox}[1]{
    colback=casegrey,
    colframe=matchaiblue,
    fonttitle=\bfseries,
    title=#1,
    arc=2mm,
    boxrule=1pt,
    enhanced,
    drop shadow
}

% =========================================================
% 3. CẤU HÌNH MÃ NGUỒN (LISTINGS)
% =========================================================
\usepackage{listings}
\lstset{
    language=Python,
    basicstyle=\ttfamily\small,
    keywordstyle=\color{blue}\bfseries,
    commentstyle=\color{green!60!black}\itshape,
    stringstyle=\color{orange},
    numbers=left,
    numberstyle=\tiny\color{gray},
    stepnumber=1,
    numbersep=10pt,
    backgroundcolor=\color{white},
    showspaces=false,
    showstringspaces=false,
    showtabs=false,
    frame=single,
    rulecolor=\color{black!30},
    tabsize=4,
    captionpos=b,
    breaklines=true,
    breakatwhitespace=false,
    title=\lstname,
    columns=flexible,
    keepspaces=true
}

% =========================================================
% 4. TRANG BÌA (TITLE PAGE)
% =========================================================
\begin{document}

\begin{titlepage}
    \centering
    \vspace*{1cm}
    
    {\Large \textbf{TRƯỜNG ĐẠI HỌC CÔNG NGHỆ}}\\[0.5cm]
    {\large KHOA CÔNG NGHỆ THÔNG TIN}\\[2cm]
    
    \includegraphics[width=0.3\textwidth]{logo.png}\\[1.5cm] % Placeholder cho logo
    
    {\Huge \textbf{CHUYÊN KHẢO NGHIÊN CỨU KHOA HỌC}}\\[1cm]
    
    {\Large \textbf{TỐI ƯU HÓA ĐỐI CHIẾU NHÂN SỰ DỰA TRÊN CHẤM ĐIỂM LAI VÀ PHÂN TÍCH NGỮ NGHĨA ĐA NGÔN NGỮ}}\\[0.5cm]
    {\large \textit{Hệ thống MatchAI Pro: Tiếp cận dựa trên Sentence-BERT và Logic Kiểm chứng Kỹ năng}}\\[2cm]
    
    \begin{flushright}
        \textbf{Nhóm nghiên cứu:} MatchAI Tech Team\\
        \textbf{Giảng viên hướng dẫn:} TS. Nguyễn Văn X\\
        \textbf{Lĩnh vực:} Natural Language Processing (NLP)
    \end{flushright}
    
    \vfill
    {\large Hà Nội, Tháng 04 Năm 2026}
\end{titlepage}

% =========================================================
% 5. TÓM TẮT & MỤC LỤC
% =========================================================
\pagenumbering{roman}

\begin{abstract}
\addcontentsline{toc}{chapter}{Tóm tắt}
Bài nghiên cứu này trình bày về \textbf{MatchAI Pro}, một hệ thống tiên tiến sử dụng Trí tuệ nhân tạo (AI) thuật toán nhằm giải quyết nghịch lý trong tuyển dụng: sự quá tải hồ sơ kết hợp với sự mất mát ứng viên tài năng do các bộ lọc từ khóa thô sơ. Chúng tôi đề xuất mô hình lai (Hybrid Model) kết hợp giữa \textbf{Sentence-BERT (E5-Small)} để trích xuất đặc trưng ngữ nghĩa sâu và \textbf{Logic Kiểm chứng Kỹ năng (Keyword Validation)} để đảm bảo tính khắt khe của quy trình tuyển dụng.

Đóng góp chính của nghiên cứu bao gồm: (1) Xây dựng hàm chấm điểm lai dựa trên không gian Hilbert 384 chiều; (2) Kỹ thuật \textbf{Hard Negative Mining} giúp AI phân biệt các CV "treo đầu dê bán thịt chó"; (3) Cơ chế giải thích AI (\textit{Critique Engine}) giúp HR ra quyết định minh bạch. Kết quả thực nghiệm cho thấy hệ thống đạt độ tương quan Spearman 0.92, vượt xa các nền tảng PTS hiện nay như LinkedIn hay TopCV.
\end{abstract}

\newpage
\tableofcontents
\newpage
\listoffigures
\listoftables

\newpage
\pagenumbering{arabic}

% =========================================================
% 6. NỘI DUNG CHÍNH (CHAPTERS)
% =========================================================

\chapter{GIỚI THIỆU (INTRODUCTION)}
\label{ch:intro}

\section{Bối cảnh và Động lực Nghiên cứu}
Trong kỷ nguyên số hóa, thị trường tuyển dụng đang chứng kiến một sự bùng nổ về mặt dữ liệu. Theo số lượng thống kê từ các nền tảng tuyển dụng lớn, một doanh nghiệp tầm trung có thể nhận tới hơn $10,000$ hồ sơ mỗi tháng cho các vị trí phổ thông. Sự quá tải này dẫn đến hiện tượng "nghẽn cổ chai" trong quy trình sàng lọc hồ sơ ứng viên.

\section{Phát biểu Bài toán}
Mặc dù các hệ thống ATS (Applicant Tracking System) đã ra đời, nhưng chúng chủ yếu dựa trên kỹ thuật \textit{Boolean Search} và \textit{Keyword Matching}. Điều này dẫn đến hai sai lầm nghiêm trọng:
\begin{itemize}
    \item \textbf{False Negative}: Loại bỏ những ứng viên xuất sắc chỉ vì họ sử dụng thuật ngữ đồng nghĩa nhưng khác với JD.
    \item \textbf{False Positive}: Chấp nhận những CV được tối ưu hóa từ khóa nghệ thuật nhưng thiếu năng lực thực tiễn.
\end{itemize}

\section{Mục tiêu Nghiên cứu}
Nghiên cứu này nhằm mục tiêu xây dựng một framework \textbf{MatchAI Pro} có khả năng mô phỏng tư duy của một chuyên gia tuyển dụng cấp cao. Chúng tôi hướng tới việc hiểu ngữ nghĩa của ứng viên thay vì chỉ nhìn vào mặt chữ.

\chapter{NGHIÊN CỨU LIÊN QUAN & SO SÁNH CẠNH TRANH}
\label{ch:related}

\section{Phân tích Đối thủ Cạnh tranh}
Trong phần này, chúng ta sẽ thực hiện đối bảng so khớp trực diện giữa MatchAI Pro và các nền tảng phổ biến như LinkedIn hay TopCV. Việc so sánh dựa trên các trụ cột công nghệ được mô tả trong \cref{tab:competitor}.

\begin{table}[H]
\centering
\caption{So sánh đặc tính kỹ thuật giữa MatchAI Pro và các nền tảng truyền thống}
\label{tab:competitor}
\begin{tabular}{@{}lll@{}}
\toprule
\textbf{Tiêu chí} & \textbf{LinkedIn/TopCV} & \textbf{MatchAI Pro} \\ \midrule
Cơ chế cốt lõi & Keyword-based & Semantic (Sentence-BERT) \\
Xử lý đa ngôn ngữ & Thường yêu cầu dịch thuật & Cross-lingual Mapping (E5) \\
Độ nhạy chuyên môn & Dễ bị đánh lừa bởi buzzwords & Hybrid Logic \& Penalty System \\
Tính giải thích & Chỉ trả ra điểm số \% & AI Critique (Giải thích lý do) \\
Quyền riêng tư & Cloud-dependent & Local Deployment Option \\ \bottomrule
\end{tabular}
\end{table}

\chapter{CƠ SỞ LÝ THUYẾT VÀ TOÁN HỌC}
\label{ch:theory}

\section{Cơ chế Self-Attention trong Transformer}
Mô hình E5 mà chúng tôi sử dụng được kế thừa từ kiến trúc Transformer. Hàm Attention được định lý hóa như sau:
\begin{equation}
    Attention(Q, K, V) = \text{softmax}\left(\frac{QK^T}{\sqrt{d_k}}\right)V
\end{equation}
Trong đó $Q, K, V$ tương ứng với Query, Key, và Value.

\section{Hàm mất mát Multiple Negatives Ranking Loss (MNRL)}
Để đạt được độ chính xác cao trong bài toán Retrieval (Tìm kiếm), chúng tôi tối ưu mô hình thông qua hàm mất mát MNRL với tham số nhiệt độ $\tau$:
\begin{equation}
\label{eq:mnrl}
L = -\sum_{i=1}^{n} \log \frac{e^{sim(q_i, p_i) / \tau}}{\sum_{j=1}^{n} e^{sim(q_i, p_j) / \tau}}
\end{equation}
Công thức này ép các vector biểu diễn của các cặp JD-CV phù hợp ($q_i, p_i$) xích lại gần nhau, đồng thời đẩy xa các cặp không phù hợp ($q_i, p_j$).

\chapter{PHƯƠNG PHÁP NGHIÊN CỨU (METHODOLOGY)}
\label{ch:methods}

\section{Kiến trúc Hệ thống Hybrid Scoring Pro}
Chúng tôi đề xuất một mô hình toán học tổng quát để tính toán mức độ phù hợp cuối cùng:
\begin{equation}
    S_{final} = \alpha \cdot \text{Sim}_{sem} + (1-\alpha) \cdot \text{Sim}_{kw} - \sum_{i \in \text{Missing}} Penalty_i
\end{equation}
Với $\alpha=0.7$, tỷ lệ này giúp hệ thống giữ được sự linh hoạt về mặt ngôn từ nhưng vẫn đảm bảo tính khắt khe về mặt kỹ thuật.

\section{Biểu diễn Logic qua Mã giả (Pseudocode)}
Để làm rõ cơ chế hoạt động, chúng tôi trình bày đoạn logic nòng cốt của server backend trong \cref{lst:hybrid_logic}.

\begin{lstlisting}[caption={Thuật toán Hybrid Scoring thực thi trong MatchAI Pro}, label={lst:hybrid_logic}]
def calculate_professional_score(cv_data, jd_data):
    # 1. Encoding van ban sang khong gian vector 384D
    v_semantic = model.encode(f"passage: {cv_data.text}")
    sim_sem = cosine_similarity(v_semantic, jd_data.vector)
    
    # 2. Calibration: Chuan hoa sim sang thang diem 0-100
    norm_sem = max(0, min(1, (sim_sem - 0.75) / 0.20))
    
    # 3. Keyword Match (Sự thật kỹ thuật)
    matched = jd_data.skills.intersection(cv_data.skills)
    kw_score = len(matched) / len(jd_data.skills)
    
    # 4. Hybrid Combination & Penalty Logic
    final_score = (norm_sem * 0.7 + kw_score * 0.3) * 100
    for skill in jd_data.critical_skills:
        if skill not in cv_data.skills:
            final_score -= 8.0 # Critical skill penalty
            
    return round(final_score, 1)
\end{lstlisting}

\chapter{THỰC NGHIỆM VÀ CASE STUDY (PHÂN TÍCH ĐỊNH TÍNH)}
\label{ch:experiments}

\section{Kịch bản Đối chứng Case Study}
Chúng tôi thực hiện phân tích trên vị trí \textbf{Senior Backend Engineer} (yêu cầu: Python, SQL, AWS, Docker).

\begin{casebox}{Case Study 1: Ứng viên Nguyễn Văn A (Positive)}
    Nguyễn Văn A sở hữu CV với đầy đủ các stack yêu cầu. Mô hình AI nhận diện được ngữ cảnh mentoring và senior-level qua các đoạn văn bản.
    \begin{itemize}
        \item Semantic similarity: 0.92
        \item Keywords matched: 7/7
        \item \textbf{Kết quả}: 94.5\% - EXCELLENT
    \end{itemize}
\end{casebox}

\begin{casebox}{Case Study 2: Ứng viên Lê Văn B (Hard Negative)}
    Lê Văn B là lập trình viên PHP, có đề cập đến "Cloud" và "Docker" dưới dạng sở thích/tìm hiểu.
    \begin{itemize}
        \item Semantic similarity: 0.81 (vẫn cao do cùng ngành IT).
        \item Keywords matched: 0/4 kỹ năng cốt lõi bắt buộc của JD.
        \item \textbf{Kết quả}: 32.8\% - REJECTED (Do hệ thống Penalty vạch trần).
    \end{itemize}
\end{casebox}

\section{Hình ảnh Minh họa}
Dưới đây là sơ đồ pipeline xử lý của hệ thống.

\begin{figure}[H]
    \centering
    \includegraphics[width=0.8\textwidth]{pipeline.png} % Thay file pipeline.png vao day
    \caption{Sơ đồ luồng xử lý dữ liệu từ PDF đến Dashboard kết quả}
    \label{fig:pipeline}
\end{figure}

\chapter{KẾT LUẬN VÀ HƯỚNG PHÁT TRIỂN}
\label{ch:conclusion}

\section{Đóng góp của Nghiên cứu}
Nghiên cứu này đã thành công trong việc kết hợp giữa AI hiện đại và logic nghiệp vụ. Hệ thống MatchAI Pro chứng minh rằng việc tin tưởng hoàn toàn vào AI là mạo hiểm, nhưng việc kết hợp nó với "Sự thật từ khóa" sẽ tạo ra một bộ lọc hoàn hảo.

\section{Hướng phát triển trong tương lai}
Chúng tôi dự kiến tích hợp cơ chế \textbf{Agentic AI} để tự động hóa hoàn toàn quy trình hội thoại ban đầu với ứng viên.

% =========================================================
% 7. TÀI LIỆU THAM KHẢO & PHỤ LỤC
% =========================================================

\begin{thebibliography}{9}
\addcontentsline{toc}{chapter}{Tài liệu tham khảo}
\bibitem{sbert} Reimers, N., \& Gurevych, I. (2019). \textit{Sentence-BERT: Sentence Embeddings using Siamese BERT-Networks}. EMNLP.
\bibitem{e5} Wang, L., et al. (2024). \textit{Multilingual E5 Text Embeddings: A Technical Report}. arXiv:2402.05672.
\bibitem{attention} Vaswani, A., et al. (2017). \textit{Attention is All You Need}. NeurIPS.
\end{thebibliography}

\appendix
\chapter{MÃ NGUỒN backend (TRÍCH ĐOẠN)}
Toàn bộ mã nguồn backend có thể được tìm thấy trong file \texttt{api\_server.py}. Dưới đây là trích đoạn cấu hình model.

\begin{lstlisting}
# GiaiDoan1_Preprocess_Code/api_server.py
try:
    model = SentenceTransformer(MODEL_PATH, device="cpu")
    print("Loaded fine-tuned model on CPU")
except:
    model = SentenceTransformer("intfloat/multilingual-e5-small")
\end{lstlisting}

\end{document}
